% !TeX root = ../papers/nc-p3.tex
\section{階数十六の四窓から二文字ブロックを回復する}
\label{sec:rank-two-block-recovery}
\begin{theorem}\label{thm:rank-two-block-recovery}
定理\ref{thm:kummer-ramification}の有理重根点と、
定理\ref{thm:klein-elliptic-completion}の$(z,\lambda)=(-2,-2)$の有理単根点を考える。
それぞれを含む印付き族の開近傍で、四辺の各関係五平面の階数二以下の曲線は
次数4と16の二つの幾何学的に整な成分を持つ。
各辺だけから両端の四文字空間の二つの二次元ブロックを、共通の交換を除き一意に回復できる。
第一辺からはさらに$z+1/z$を回復できる。これらは添字保持同型の必要不変量である。
\end{theorem}
\begin{proof}
\path{computations/verify_rank_two_block_recovery.py}は各標本の二切断と
それらの転置について、前節の有限超平面切断を同一素数1009で取り直す。
八つの曲線それぞれで残る成分軌道は$20$か$4+16$である。
厚い二色の四平面上で対角二ブロックの行列式が切る四次曲線があるため、全て後者となる。
奇位置で矢の順を交換しても行列の転置なので成分次数は変わらない。

開近傍への移行を詳しく述べる。期待次元の線型切断と厚い二行列式の
固有交叉を保つ開部分では、その次数20と4の交叉サイクルは族をなす。
これは固定次数の固有交叉のChow形式、または対応するCohen--Macaulayな
定次元の切断の平坦性による。幾何学的に二つ以上の有効サイクルの和に分解することは
次数ごとのChow多様体の加法の固有性から閉条件である。
まず四次の分解を除く。次に全サイクルが指定四次と、和の次数が16となる
二つの正次数サイクルに分解する閉条件を全て除く。
良い還元の二つの整成分はどの条件にも属さないので、標数零の二標本も残る。
これで所要の幾何学的整成分と重複度が開近傍で確定する。

次数4成分は厚い四次元関係空間$D$の射影内の二次式完全交叉であり、
一次式はそのイデアルにないのでその線型包が$\mathbb P(D)$となる。
辺の行列を第二矢空間の双対から第一矢空間への写像とみなす。
$SM=MT$が全ての$M\in D$について成立する自己準同型の対$(S,T)$を求める。
四つの行列についての64本の線型式は32未知数に対し、八辺全ての有理標本で階数30である。
二つの既知ブロックの同時射影が解なので、階数30を保つ開集合で解代数は$k\oplus k$である。
順序を付けない二つの原始べき等元はこの代数だけで決まる。それらの像および
双対像により両矢の二ブロックを取り戻せる。二変数形式の商に必要な次数5の
階数二曲線もどの辺にもない。

第一辺で、回復した二べき等元を用いて対角ブロックを除くと薄い二色の鉛筆を得る。
その二つの角ブロックの行列式は$r^2-t^2,zr^2-t^2$である。
従って鉛筆上の二つの零点対は相異なる四点で、二対の中を同時に交換する
射影直線の自己同型は一意である。現在の表示では$t/r\mapsto-t/r$であり、
その二固定点が薄い二色の方向を与える。二固定点を零と無限遠に取り、
別々の零点対から一つずつ選ぶと座標比の平方は$z$または$1/z$である。
座標の倍率、各対内の選択、二つの固定点の順序、二ブロックの交換によっても
$z+1/z$は変わらない。この計算は四元群の印を入力として必要としない。
\end{proof}

\begin{corollary}\label{cor:three-character-frames}
上の二標本を含む開集合を縮めれば、どの辺にも、各二文字ブロックを二本の
文字直線に分ける候補は三通りで足りる。四位置を通じても、互いに整合する
順序なし文字直線の分解は高々三通りである。
\end{corollary}
\begin{proof}
$D$の低・高ブロックの行列式の対称行列を$O,P$とする。いずれも非退化である。
コードで八辺の全てについて$O^{-1}P$の固有値が相異なり、トレースが非零であることを
有理数上で確認する。これらは開条件であり、両行列式の単元倍にも依存しない。
二つの二次曲面をそれぞれ保つ射影変換は$O^{-1}P$をスカラー倍まで保つ。
トレースからその倍率は1である。よって四固有直線を個々に固定し、
$O$について直交する固有基底で各係数は共通倍率を除き符号である。
両二次曲面は二行二列の階数一行列のSegre曲面であり、両方の母線族を保つには
符号の積が正でなければならない。実際、直交四次元空間で一座標だけの符号反転は
極大等方二平面をそれと直線で交わる平面に写し、二族を交換する。
従って残る群は同時符号を法とした偶符号の四元群である。

非自明三元の各々の二つの固有二平面では両形式とも非退化である。
ゆえに各Segre曲面の固定点は四点で、両方の射影直線因子で非自明な対合として作用する。
どちらの曲面の左・右因子でも、この四元群は忠実である。非自明な二対合の固定点対が
交わることはない。実際、一点を固定する射影直線の有限群は、その点を無限遠にして
有限軌道の平均を原点とすれば乗法群に入り、巡回群である。
もとの文字の下位ビットの符号は、低・高の各ブロックの各矢で非自明で、
厚い関係に作用し二つの行列式と各母線族を保つ。従ってこの三元の一つであり、
四つの二文字空間内の固定点対が求める文字直線を含む候補を与える。
隣の辺では一つの共通ブロックの固定点対を指定すれば候補は高々一つなので、
一周しても数は増えない。これは候補の上界であり、薄い係数の照合でさらに減る場合を含む。
\end{proof}
